% TOP_PROBLEM: {"id": 3138, "title": "Cycle Double Covers Containing Predefined 2-Regular Subgraphs", "category_id": 3, "category": {"id": 3, "name": "graph_theory", "display_name": "Graph Theory", "description": "Problems involving graphs, networks, and their properties.", "slug": "graph-theory", "order_index": 3, "created_at": "2026-07-31T15:26:25.671Z"}, "status": "open", "problem_number": "OPG-60029", "set_id": 12}
% FIRST_POSTED: 2026-10-01
% ATTEMPT_STATUS: unresolved
% UnsolvedMath ID 3138; source catalogue code OPG-60029
% !TeX program = lualatex
% Research attempt by GPT-6 Astra Ultra; no claim of a new complete solution.
\documentclass[11pt,a4paper]{article}
\usepackage[margin=25mm]{geometry}
\usepackage{fontspec}
\setmainfont{lmroman10-regular.otf}[BoldFont=lmroman10-bold.otf,
 ItalicFont=lmroman10-italic.otf,BoldItalicFont=lmroman10-bolditalic.otf]
\usepackage{amsmath,amssymb,mathtools}
\usepackage{unicode-math}
\setmathfont{latinmodern-math.otf}
\AtBeginDocument{\let\setminus\smallsetminus}
\usepackage{xurl,hyperref}
\hypersetup{colorlinks=true,urlcolor=blue}
\setcounter{secnumdepth}{0}
\setlength{\parindent}{0pt}
\setlength{\parskip}{5pt}
\setlength{\emergencystretch}{3em}
\begin{document}
\section*{Cycle Double Covers Containing Predefined 2-Regular Subgraphs}\label{3138}
{\small UnsolvedMath ID 3138 · OPG-60029 · Graph Theory · OpenGarden}
\subsection{Definitions and mathematical statement}
Let $G=(V,E)$ be a finite simple cubic graph that is $2$-connected:
it has at least three vertices and remains connected after deletion
of any one vertex. A circuit means a connected $2$-regular subgraph,
or equivalently a simple cycle. Let $S$ be a $2$-regular subgraph
of $G$, possibly with several components, and suppose that the
spanning graph $(V,E\setminus E(S))$ is connected. Deleting the edges
of $S$ retains all vertices of $G$.

A cycle double cover of $G$ is a finite list $\mathcal C$ of circuits
such that, counting appearances in the list,
\[
 \bigl|\{C\in\mathcal C:e\in E(C)\}\bigr|=2
 \qquad(e\in E).
\]
Repetitions are permitted by the general cover convention. To say
that the cover contains $S$ means that every circuit component of
$S$ occurs as an entire member of $\mathcal C$. Merely covering the
edges of $S$ within other circuits does not meet this requirement.

The conjecture asserts that such a cover containing $S$ exists for
every pair $(G,S)$ satisfying the hypotheses. There is no prescribed
bound on the number of covering circuits and no requirement that
they be disjoint. The connectivity assumption is imposed on the
edge complement of the specified subgraph, and $S$ is not assumed
to be spanning. These distinctions determine which prescribed
cycle families the extension statement addresses.
\subsection{Short English statement}
Can every prescribed union of disjoint cycles whose removal leaves a cubic graph connected be retained inside a cycle double cover?
\subsection{Research attempt}
\paragraph{Scope and process.}
This is a GPT-6 Astra Ultra research attempt dated 1 October 2026, using
primary-source browsing and elementary mathematical checks. The canonical
definition above is retained. Results below are restricted results or
reductions, not a claim of a new solution to the entire catalogue question.
No formal proof assistant or external mathematical referee checked this text.


\paragraph{Literature and conventions.}
Hoffmann-Ostenhof, Zhang and Zhang [R1] study prescribed $2$-regular
subgraphs and prove sufficient conditions, including particular
tree-and-cycle decompositions. The target here retains every
component of $S$ as an entire circuit. We prove two restricted
extension mechanisms and explain why parity alone cannot replace
that requirement. In this text a circuit is simple and connected,
exactly as in the retained statement.

\paragraph{A local necessary transition rule.}
At a vertex of a cubic graph, name the incident edges $a,b,c$.
Every circuit passing through the vertex uses exactly one pair
of these edges. Let $x_{ab},x_{ac},x_{bc}$ count the occurrences
of the three pairs in a cycle double cover. Covering every edge
twice requires
\[
 x_{ab}+x_{ac}=2,\quad x_{ab}+x_{bc}=2,\quad x_{ac}+x_{bc}=2.
\]
The unique solution is $x_{ab}=x_{ac}=x_{bc}=1$. Thus each local
pair must be used exactly once. A specified circuit fixes one
of these transitions at every vertex it visits; a proof must
complete the other transitions coherently around whole circuits.
Counting even edge multiplicities alone omits this compatibility.

\paragraph{A full extension in a compatible edge-colouring subcase.}
Suppose $G$ has a proper edge-colouring by $1,2,3$, and every
component of $S$ is bichromatic in this colouring. For each pair
$\{i,j\}$ take the subgraph consisting of edges of those two
colours. Every vertex has one edge of each colour, so this is
a spanning $2$-regular subgraph; split it into its simple circuit
components. Collect the components from all three colour pairs.
An edge of colour $i$ appears in exactly the two pairs containing
$i$, so the resulting list is a cycle double cover.

If a component $K$ of $S$ uses colours $i,j$, then at every vertex
of $K$ its two edges already exhaust the two edges in that
colour-pair subgraph. It follows that $K$ is an entire component
of that subgraph and hence occurs in the cover as required.
Different components of $S$ may use different pairs. This argument
does not need the edge-complement connectivity hypothesis, but it
does need the extra compatible-colouring hypothesis.

\paragraph{A separate planar subcase.}
Suppose a $2$-connected plane embedding of $G$ is given and every
component of $S$ is the boundary of a face. In a $2$-connected
plane graph each facial boundary is a simple circuit: a repeated
boundary vertex would separate the portions encountered between
two visits and be a cut vertex. Every edge has two sides, and
the facial boundary circuits, including that of the unbounded
face, therefore cover every edge exactly twice. They provide
the desired cover containing $S$.

These two mechanisms are not interchangeable. In $K_4$, a triangle
$S$ has connected edge complement (a three-edge star). The four
triangles are a double cover and one may embed any chosen triangle
as a face. But a triangle cannot be bichromatic in a proper
edge-colouring: a two-coloured circuit must alternate colours and
therefore have even length. The first mechanism misses this valid
instance while the facial construction handles it.

\paragraph{A tempting Eulerian shortcut and its precise defect.}
Replace each edge of $G$ by two parallel copies and remove one
copy of each edge in $S$. The remaining multigraph has degree four
at vertices of $S$ and degree six elsewhere, so it decomposes
into circuits in the multigraph sense. Adding back $S$ then has
the correct edge multiplicities. However, those multigraph circuits
may be two-edge circuits formed by the parallel copies of a
single original edge. Their projection to $G$ is not a simple
circuit. Choosing such pairs for all doubled edges is precisely
why even degree by itself gives no cycle double cover of an
arbitrary bridgeless graph. A proof must exclude these projected
backtracks and preserve the prescribed circuits simultaneously.

\paragraph{Verification and first remaining gap.}
The edge-colouring proof tracks individual circuit components,
not just the unions of two colour classes. The plane argument
requires a compatible embedding; connectivity of $G-E(S)$ does
not itself provide one or make $G$ planar. The local transition
equations are necessary but contain no global assurance of simple
projected circuits. \textbf{UNRESOLVED}: the remaining step is a
general extension construction from the stated connectivity
hypothesis alone, without either additional colouring or embedding.

\subsection{Sources}
\begin{itemize}
\item[S1] ULAM.AI and the UnsolvedMath Contributors, supplied problems.json, record 3138 (OPG-60029), OpenGarden collection. Catalogue identity and source transcription; CC BY 4.0.
\url{https://huggingface.co/datasets/ulamai/UnsolvedMath}.
\item[S2] Open Problem Garden, Cycle Double Covers Containing Predefined 2-Regular Subgraphs. Original problem and discussion; the mathematical formulation below is expanded from the supplied source record.
\url{http://www.openproblemgarden.org/op/cycle_double_covers_containing_predefined_2_regular_subgraphs}.

\item[R1] Arthur Hoffmann-Ostenhof, Cun-Quan Zhang and Zhang Zhang,
\emph{Cycle double covers and non-separating cycles}, European
Journal of Combinatorics 81 (2019), 276--284; primary preprint.
\url{https://arxiv.org/abs/1711.10614}.

\end{itemize}
\end{document}
